\begin{proofbox}
\textbf{\textcolor{CProof}{思路提示}}
先独立证明核心部分 $GM \le AM$（基本不等式），再利用倒数关系和平方关系，依次向两端（HM 和 QM）延伸，最终串联成链。

\medskip
\textbf{\textcolor{CProof}{正式推导}}
\begin{description}[style=nextline, leftmargin=2.8em, labelsep=0.5em, itemsep=0.55em]
    \item[\textbf{步骤一（证明 GM≤AM）：}] 由基本不等式 $(\sqrt{a} - \sqrt{b})^2 \ge 0$ 展开可得。
    \[
    \begin{aligned}
    (\sqrt{a} - \sqrt{b})^2 &\ge 0 \\
    a - 2\sqrt{ab} + b &\ge 0 \\
    a + b &\ge 2\sqrt{ab} \\
    \frac{a+b}{2} &\ge \sqrt{ab}.
    \end{aligned}
    \]
    \par\textbf{理由：}完全平方的非负性是基本不等式的来源。等号成立当且仅当 $\sqrt{a} = \sqrt{b}$，即 $a=b$。

    \item[\textbf{步骤二（证明 HM≤GM）：}] 对上一步结论 $\frac{a+b}{2} \ge \sqrt{ab}$ 中的 $a, b$ 同时取倒数，并利用不等式方向不变的性质。
    \[
    \begin{aligned}
    \text{令 } a' = \frac{1}{a},\; b' = \frac{1}{b} &\quad (a', b' > 0) \\
    \frac{a'+b'}{2} &\ge \sqrt{a'b'} \\
    \frac{\frac{1}{a}+\frac{1}{b}}{2} &\ge \sqrt{\frac{1}{ab}} = \frac{1}{\sqrt{ab}} \\
    \frac{1}{a}+\frac{1}{b} &\ge \frac{2}{\sqrt{ab}} \\
    \frac{2}{\frac{1}{a}+\frac{1}{b}} &\le \sqrt{ab}.
    \end{aligned}
    \]
    \par\textbf{理由：}正实数的倒数仍是正实数，不等式两边同时取倒数（并注意不等号方向因正数而反向）。

    \item[\textbf{步骤三（证明 AM≤QM）：}] 再次利用平方非负性 $(a-b)^2 \ge 0$ 的变形。
    \[
    \begin{aligned}
    (a-b)^2 &\ge 0 \\
    a^2 - 2ab + b^2 &\ge 0 \\
    a^2 + b^2 &\ge 2ab \\
    2(a^2 + b^2) &\ge a^2 + 2ab + b^2 \\
    2(a^2 + b^2) &\ge (a+b)^2 \\
    \frac{(a+b)^2}{4} &\le \frac{a^2+b^2}{2} \\
    \frac{a+b}{2} &\le \sqrt{\frac{a^2+b^2}{2}}.
    \end{aligned}
    \]
    \par\textbf{理由：}从 $(a-b)^2 \ge 0$ 出发，逐步配凑出 $(a+b)^2$ 与 $a^2+b^2$ 的关系，最后开方（注意两边均为正，开方保序）。
\end{description}

\medskip
\textbf{\textcolor{CProof}{结论回扣}}
\textbf{综合步骤一、二、三，我们得到了完整的不等式链：}
$\frac{2}{\frac{1}{a}+\frac{1}{b}} \le \sqrt{ab} \le \frac{a+b}{2} \le \sqrt{\frac{a^{2}+b^{2}}{2}}$。\textbf{每一步的取等条件均为 $a=b$，故整个不等式链等号成立当且仅当 $a=b$。}
\end{proofbox}
